\textit{Proof.}
Put
\[
 a:=\frac1n,\qquad v:=\frac d{n^2},\qquad
 m:=\min(u,s).
\tag{1}
\]
In the nonsaturated triangular regime \(L^2<d<nL\), one has
\[
 0<v<u<1,
\tag{2}
\]
because \(u/v=d/L^2>1\) and \(u=(d/(nL))^2<1\).  Let \(t:=\min\{u,s+v\}\).  Since \(v<u\) and \(v\le s+v\), we have \(t\ge v\).  Also \(t\ge m\): if \(s\le u\), then both arguments in the minimum defining \(t\) are at least \(s=m\), while if \(u<s\), then \(t=u=m\).  Consequently
\[
 q=a+t\ge a+\max(v,m)
 \ge\frac12(a+v+m)=\frac12\{b+\min(u,s)\}.
\tag{3}
\]
For the reverse inequality, the elementary inequality
\[
 \min\{u,s+v\}\le v+\min(u,s)
\tag{4}
\]
follows by separating the cases \(u\le s\) and \(s<u\).  Thus
\[
 q\le a+v+m=b+\min(u,s).
\tag{5}
\]
Because \(u<1\), the normalized converse benchmark is exactly
\[
 \ell=b+s\min(1,u)=b+su.
\tag{6}
\]
This proves the first display.

We next compare the benchmarks globally.  Set
\[
 x:=\min(1,u),\qquad
 r:=a+\min\{x,s+v\},\qquad
 \Lambda:=b+sx.
\tag{7}
\]
Theorems~\(\mathrm{thm:frontier\text{-}upper}\) and \(\mathrm{thm:radius\text{-}channel\text{-}converse}\) give constants depending only on \(\epsilon\) for which
\[
 c_\epsilon M^2\Lambda
 \le\mathsf R_{n,d,\epsilon,M,\sigma}
 \le C_\epsilon M^2r.
\tag{8}
\]
If \(\sigma\ge\underline\sigma>0\), then \(s\ge\underline\sigma^2\), and
\[
 \Lambda\ge a+\underline\sigma^2x
 \ge\min(1,\underline\sigma^2)(a+x)
 \ge\min(1,\underline\sigma^2)r.
\tag{9}
\]
Combining (8) and (9) proves the fixed-positive-radius equivalence.

Now fix \(C<\infty\).  If \(u\ge1\), then \(x=1\) and
\[
 r=a+\min(1,s+v)\le a+s+v=\Lambda.
\tag{10}
\]
If \(u\le Cb\), then \(x\le u\le Cb\), so
\[
 r\le a+x\le(1+C)b\le(1+C)\Lambda.
\tag{11}
\]
If \(s\le Cb\), then
\[
 r\le a+s+v=b+s\le(1+C)b\le(1+C)\Lambda.
\tag{12}
\]
Equations (8), (10)--(12) give the claimed order equivalence with constants uniform in \(\sigma\) on each displayed regime.

The condition \(u\le b\) is equivalent, after multiplication by \(n^2L^2>0\), to
\[
 d^2\le nL^2+dL^2.
\tag{13}
\]
If \(d\le\sqrt n\,L\), then \(d^2\le nL^2\), so (13) holds.  If \(d\le L^2\), then \(d^2\le dL^2\), so (13) again holds.  Finally, \(s\le Cb\) is precisely the stated radius-dominance elbow.

It remains to localize possible divergence.  If \(x\ge\eta>0\), then
\[
 r\le b+s,
 \qquad
 \Lambda=b+sx\ge\min(1,\eta)(b+s),
\tag{14}
\]
so the benchmark ratio is bounded.  Equations (9)--(14) show that a sequence along which the upper benchmark divided by the lower benchmark diverges can have none of the following properties along an infinite subsequence: \(s\) bounded away from zero, \(x\) bounded away from zero, \(u\ge1\), \(u=O(b)\), or \(s=O(b)\).  Therefore such a sequence must satisfy
\[
 b\ll u\ll1,
 \qquad b\ll s\ll1.
\tag{15}
\]
Conversely, (15) is the only regime not covered by the preceding constant-factor comparisons; this does not assert that divergence occurs at every point of the wedge.  Since \(L=o(\sqrt n)\), the relations \(u\gg b\) and \(u\ll1\) are, up to boundary constants, equivalent to
\[
 \sqrt n\,L\ll d\ll nL.
\tag{16}
\]
Indeed, \(u/(1/n)=d^2/(nL^2)\), while \(u/(d/n^2)=d/L^2\), and the first divergence in (16) implies the second.

For the exhibited sequence \(d_n=n\), \(L_n=\log(en)\), and \(\sigma_n=L_n^{-1/2}\),
\[
 b_n=\frac2n,\qquad u_n=\frac1{L_n^2},\qquad s_n=\frac1{L_n}.
\tag{17}
\]
Hence the selector benchmark is
\[
 \frac1n+\min\left\{1,\frac1{L_n^2},\frac1{L_n}+\frac1n\right\}
 =\frac1n+\frac1{L_n^2}
 \asymp\frac1{L_n^2},
\tag{18}
\]
whereas the proved lower benchmark is
\[
 \frac2n+\frac1{L_n^3}
 \asymp\frac1{L_n^3}.
\tag{19}
\]
Their ratio is therefore of order \(L_n\to\infty\).

The endpoint \(\sigma=0\) is covered by (12), every fixed positive radius by (9), saturation by (10), the small-alphabet and first parametric elbows by (11), and the shrinking-radius parametric elbow by (12).  Thus (15) is exactly the residual regime not decided by the established selector and radius channel. \(\square\)